% Part: first-order-logic
% Chapter: introduction
% Section: satisfaction

\documentclass[../../../include/open-logic-section]{subfiles}

\begin{document}

\olfileid[hi]{fol}{int}{sat}

\olsection{संतुष्टि}

!!{formula}s की परिभाषा समझ लेने पर हम प्रथम-क्रम तर्क के अर्थविज्ञान पर
आगे बढ़ सकते हैं। यहाँ मूल परिभाषा !!a{structure} की है। हमारी सरल भाषा के
लिए !!a{structure}~$\Struct M$ के केवल तीन घटक हैं: \emph{!!{domain}} कहलाने
वाला कोई अरिक्त समुच्चय $\Domain{M}$, वह वस्तु जिसे~$\Obj a$,~$\Struct M$
में निर्दिष्ट करता है, और वे वस्तुएँ जिनके लिए~$\Obj P$,~$\Struct M$ में
सत्य है। $\Obj a$ द्वारा निर्दिष्ट वस्तु को~$\Assign{\Obj a}{M}$ और उन
वस्तुओं के समुच्चय को, जिनके लिए~$\Obj P$ सत्य है,
~$\Assign{\Obj P}{M}$ लिखा जाता है। !!^a{structure}~$\Struct{M}$ में ठीक ये
तीन चीजें होती हैं: $\Domain{M}$, $\Assign{\Obj a}{M} \in
\Domain{M}$ और $\Assign{\Obj P}{M} \subseteq \Domain{M}$। सामान्य स्थिति
अधिक जटिल होगी, क्योंकि उसमें अनेक !!{predicate}s और !!{constant}s का उपयोग होगा,
!!{constant}s के एक से अधिक स्थान हो सकते हैं, और !!{function}s का भी उपयोग होगा।

यह उन !!{formula}s के लिए संतुष्टि परिभाषित करने को पर्याप्त है जिनमें
!!{variable}s की कोई उपस्थिति नहीं है। विचार ऐसी आगमनात्मक परिभाषा देने का है जो
!!{formula}s की हमारी परिभाषा के अनुरूप हो। हम निर्दिष्ट करते हैं कि कोई
परमाण्विक सूत्र~$\Struct{M}$ में कब संतुष्ट है; फिर, उदाहरणार्थ, $\lnot !A$
~$\Struct{M}$ में कब संतुष्ट है, यह इस आधार पर बताते हैं कि $!A$,
~$\Struct{M}$ में संतुष्ट है या नहीं। हम इस प्रकार परिभाषित कर सकते हैं:
\begin{enumerate}
  \item $\Atom{\Obj P}{\Obj a}$,~$\Struct{M}$ में तभी और केवल तभी संतुष्ट
  है जब $\Assign{\Obj a}{M} \in \Assign{\Obj P}{M}$।
  \item $\lnot !A$,~$\Struct{M}$ में तभी और केवल तभी संतुष्ट है जब $!A$,
  ~$\Struct{M}$ में संतुष्ट नहीं है।
  \item $(!A \land !B)$,~$\Struct{M}$ में तभी और केवल तभी संतुष्ट है जब
  $!A$,~$\Struct{M}$ में संतुष्ट है और $!B$,~$\Struct{M}$ में भी संतुष्ट है।
\end{enumerate}
मान लें कि $\Domain{M} = \{0, 1, 2\}$, $\Assign{\Obj a}{M} = 1$ और
$\Assign{\Obj P}{M} = \{1, 2\}$। यह परिभाषा बताती है कि
$\Atom{\Obj P}{\Obj a}$,~$\Struct{M}$ में संतुष्ट है (क्योंकि
$\Assign{\Obj a}{M} =  1 \in \{1,2\} = \Assign{\Obj P}{M}$)। वह यह भी
बताती है कि $\lnot \Atom{\Obj P}{\Obj a}$,~$\Struct{M}$ में संतुष्ट नहीं
है; फलतः $\lnot\lnot \Atom{\Obj P}{\Obj
a}$ संतुष्ट है और
$(\lnot \Atom{\Obj P}{\Obj a} \land \Atom{\Obj P}{\Obj a})$ संतुष्ट नहीं
है; इत्यादि।

कठिनाई तब आती है जब हम परिमाणकों की परिभाषा देना चाहते हैं। हम कुछ ऐसा
कहना चाहेंगे: ``$\lexists[\Obj
v_0][\Atom{\Obj P}{\Obj v_0}]$ तभी और केवल तभी संतुष्ट है जब
$\Atom{\Obj P}{\Obj
v_0}$ संतुष्ट है।'' किंतु !!{structure}~$\Struct{M}$ यह नहीं बताती कि
!!{variable}s के साथ क्या किया जाए। वास्तव में हम यह कहना चाहते हैं कि
$\Atom{\Obj P}{\Obj v_0}$, \emph{$\Obj v_0$ के किसी मान के लिए} संतुष्ट है।
इसे सटीक करने के लिए हमें~$\Domain{M}$ के !!{element}s में से किसी को केवल~$\Obj a$ हेतु ही
नहीं, बल्कि~$\Obj
v_0$ हेतु भी चुनने की विधि चाहिए। इसके लिए हम !!{variable}
\emph{निर्धारण} लेते हैं। !!^a{variable} निर्धारण केवल ऐसा फलन~$s$ है जो
!!{variable}s को~$\Domain{M}$ के !!{element}s पर प्रतिचित्रित करता है (हमारे
उदाहरण में $1$, $2$ अथवा~$3$ में से किसी पर)। पहले से यह मालूम नहीं कि किसी
!!a{formula} में किन !!{variable}s की उपस्थितियाँ होंगी, इसलिए हम उन !!{variable}s को
सीमित नहीं कर सकते जिन्हें~$s$ मान देता है। सरल समाधान यह अपेक्षा करना है कि
$s$ \emph{सभी} !!{variable}s~$\Obj
v_0$, $\Obj v_1$, \dots\@ को मान दे। हम केवल आवश्यक मानों का उपयोग करेंगे।

अब हम !!{formula}s की संतुष्टि केवल !!a{structure} के सापेक्ष नहीं, बल्कि
!!a{structure}~$\Struct{M}$ \emph{और} !!a{variable} निर्धारण~$s$ के
सापेक्ष परिभाषित करेंगे और संक्षेप में $\Sat{M}{!A}[s]$ लिखेंगे। हमारी
परिभाषा में !!{variable}s वाले परमाण्विक !!{formula}s के लिए एक अतिरिक्त
खंड होगा:
\begin{enumerate}
  \item $\Sat{M}{\Atom{\Obj P}{\Obj a}}[s]$ तभी और केवल तभी जब
  $\Assign{\Obj a}{M} \in \Assign{\Obj P}{M}$।
  \item $\Sat{M}{\Atom{\Obj P}{\Obj v_i}}[s]$ तभी और केवल तभी जब
  $s(\Obj v_i) \in \Assign{\Obj P}{M}$।
  \item $\Sat{M}{\lnot !A}[s]$ तभी और केवल तभी जब $\Sat{M}{!A}[s]$ नहीं।
  \item $\Sat{M}{(!A \land !B)}[s]$ तभी और केवल तभी जब $\Sat{M}{!A}[s]$
  और $\Sat{M}{!B}[s]$ दोनों।
\end{enumerate}
इससे एक समस्या हल हो गई: अब हम कह सकते हैं कि $\Struct{M}$,
$\Atom{\Obj P}{\Obj v_0}$ को मान~$s(\Obj v_0)$ के लिए कब संतुष्ट करती है।
$\lexists[\Obj v_0][\Atom{\Obj P}{\Obj
v_0}]$ की परिभाषा ठीक करने के लिए एक काम और करना होगा। हम चाहते हैं कि
$\Sat{M}{\lexists[\Obj v_0][\Atom{\Obj P}{\Obj v_0}]}[s]$ तभी और केवल तभी
हो जब $\Sat{M}{\Atom{\Obj P}{\Obj v_0}}[s']$,~\emph{किसी} ऐसे निर्धारण
$s'$ के लिए सत्य हो जो~$\Obj v_0$ को कोई मान देता है। किंतु~$\Obj v_0$ को
दिया मान अनिवार्यतः $s(\Obj v_0)$ द्वारा निर्दिष्ट मान नहीं होना चाहिए। इसके
लिए एक संकेतन लें: यदि $m \in \Domain{M}$, तो $\Subst{s}{m}{\Obj v_0}$ वह
निर्धारण है जो~$s$ जैसा है (सभी !!{variable}s~$\Obj v_0$ को छोड़कर), किंतु
~$\Obj v_0$ को~$m$ देता है। अब हमारी परिभाषा यह हो सकती है:
\begin{enumerate}\setcounter{enumi}{4}
  \item $\Sat{M}{\lexists[\Obj v_i][!A]}[s]$ तभी और केवल तभी जब
  $\Sat{M}{!A}[\Subst{s}{m}{\Obj v_i}]$, किसी~$m \in \Domain{M}$ के लिए।
\end{enumerate}
क्या यह काम करता है? मान लें कि $s(\Obj v_i) = 0$, सभी~$i \in
\Nat$ के लिए। $\Sat{M}{\lexists[\Obj v_0][\Atom{\Obj P}{\Obj v_0}]}[s]$
तभी और केवल तभी जब कोई $m \in \Domain{M}$ ऐसा हो कि
$\Sat{M}{\Atom{\Obj P}{\Obj
v_0}}[\Subst{s}{m}{\Obj v_0}]$। ऐसा मान है: हम $m = 1$ अथवा $m = 2$ चुन
सकते हैं। ध्यान दें कि यह तब भी सत्य है जब मान~$s(\Obj v_0)$, जो~$\Obj v_0$
को स्वयं~$s$ द्वारा दिया गया है---इस स्थिति में $0$---काम नहीं करता। हमारे
पास $\Sat{M}{\Atom{\Obj P}{\Obj v_0}}[\Subst{s}{1}{\Obj
v_0}]$ है, किंतु $\Sat{M}{\Atom{\Obj P}{\Obj v_0}}[s]$ नहीं।

यदि यह उलझा हुआ और बोझिल लगता है, तो वास्तव में ऐसा ही है। किंतु सभी
!!{formula}s के लिए संतुष्टि की सटीक आगमनात्मक परिभाषा देने हेतु यह अतिरिक्त
जटिलता आवश्यक है; अर्थवैज्ञानिक धारणाओं को सटीक रूप से परिभाषित करने के लिए
हमें इसी प्रकार की कोई विधि चाहिए। इसे करने की दूसरी विधियाँ भी हैं, पर वे
सभी समान रूप से (अ)सुरुचिपूर्ण हैं।

\end{document}
